\documentclass[10pt,a4paper]{amsart}
\usepackage[margin=19mm]{geometry}
\usepackage{amsmath,amssymb,mathtools}
\usepackage[scr=rsfs,cal=euler]{mathalfa}
\usepackage{xfrac}
\usepackage[unicode,hidelinks,bookmarks=false]{hyperref}
\newcommand{\ie}{i.e.\ }
\newcommand{\cf}{cf.\ }
\newcommand{\resp}{respectively\ }
\newcommand{\Subsection}{\subsection}
\providecommand{\nobreakdash}{\nobreak\hbox{-}}
\setlength{\emergencystretch}{2em}
\multlinegap0pt
\usepackage{bm}
\usepackage{amssymb}
\newtheorem{lemma}{Lemma}
\newtheorem{theorem}{Theorem}
\newcommand{\A}{\boldsymbol{A}}
\newcommand{\B}{\boldsymbol{B}}
\newcommand{\D}{\boldsymbol{D}}
\newcommand{\I}{\boldsymbol{I}}
\newcommand{\J}{\boldsymbol{J}}
\newcommand{\Sm}{\boldsymbol{S}}
\newcommand{\jv}{\boldsymbol{j}}
\newcommand{\vv}{\boldsymbol{v}}
\title{Pseudo-Euclidean representations of switching classes of Johnson and Hamming graphs with minimal dimension}
\author{Hiroshi Nozaki, Masashi Shinohara, Sho Suda}
\date{Source: arXiv:2507.13592v2 (CC BY 4.0)}
\begin{document}
\maketitle

\noindent\emph{Source and licence.} Pseudo-Euclidean representations of switching classes of Johnson and Hamming graphs with minimal dimension. Version arXiv:2507.13592v2 (CC BY 4.0), distributed by arXiv under the Creative Commons Attribution 4.0 International licence, http://creativecommons.org/licenses/by/4.0/, as recorded in the arXiv licence metadata for this exact version and shown on the abstract page https://arxiv.org/abs/2507.13592v2. Full licence notice: \texttt{licenses/arxiv-2507.13592-CC-BY-4.0.txt}.\par\smallskip

\noindent\textit{Editorial scope.} Counted unit: the article's theorem thm:regular\_angle (source0.tex lines 582--592). The article's complete original proof of it is delivered with it (source0.tex lines 593--609, one \texttt{proof} environment of 17 lines). No mathematics is written, repaired or re-derived: the statement and the proof are the article's own lines, in the article's own notation, and no retained line is substituted or re-flowed.  Retained uncounted as the source-local context the proof needs: lines 422--433; lines 470--480; lines 487--496; lines 506--516.  Nothing was omitted from any retained span, so the package carries no omission record.  No retained line contains a citation, so the wrapper carries no bibliography and no bibliography entry.  No retained line contains a figure environment, an image-inclusion command or drawing code, so no image is delivered and the package declares no asset dependency.  Editorial formatting only, in addition: an AMS \texttt{amsart} preamble that restates the article's own declarations and macros used by the retained lines, namely the environments \texttt{lemma}, \texttt{theorem} on separate counters, together with the standard packages those lines need.  The retained environments print in this document as Lemma 1, Theorem 1 in order of appearance, whereas the article prints the same items with the numbering of its own full document.  The complete native source of the article as distributed in the arXiv e-print for this exact version is bundled unchanged as \texttt{sources/181787/source0.tex}.
\begin{lemma} \label{lem:eigen_SD}
Let $(V,E)$ be a $k$-regular graph of order $n$ with adjacency matrix $\A$. 
 Let $\lambda_0=k,\lambda_1,\ldots,\lambda_r$ be the distinct eigenvalues of $\A$, and  $E_0,E_1,\ldots,E_r$ the corresponding eigenspaces.  
Then, the following hold. 
\begin{enumerate}
\item The eigenspaces of $\A$ 
%$\overline{\A}$, $\D$, 
and $\Sm_{\D}$ are the same. 
\item The distinct eigenvalues of $\Sm_{\D}$ for $E_i$ are 
$\tau_0=(a-b)(2k-n+1)$ and $\tau_i=(a-b)(2\lambda_i+1)$ for $i\in\{1,\ldots,r\}$. 
\end{enumerate}
\end{lemma}

\begin{lemma} \label{lem:eigen_2D'}
Let $(V,E)$ be a $k$-regular graph. The notation is defined as above. 
Let $\lambda_0=k,\lambda_i$ be the eigenvalues of $\A$, 
and $\tau_i$ those of $\Sm_{\D}$, which are given in  Lemma~\ref{lem:eigen_SD}. 
Then the eigenvalues of $2\D'-(a+b)\J$ are  
\begin{align*}%\label{eq:mu0}
\mu_0&=\tau_0-(a+b)=(a-b)(2k-n+1)-(a+b),\\ %\label{eq:mui}
\mu_i&=\tau_i-(a+b)=(a-b)(2\lambda_i+1)-(a+b) \text{ for $i\ne 0$,}
\end{align*}
 and the eigenspaces associated with $\mu_i$ are $\I_U E_i$. 
\end{lemma}

\begin{lemma} \label{lem:switch_main}
Let $(V,E)$ be a $k$-regular graph. The notation is defined as above. 
Then the following are equivalent for $i= 1,\ldots,r$.  
\begin{enumerate}
\item $\mu_i$ is not a main eigenvalue of $2\D'-(a+b)\J$. 
\item $\vv^\top(\I_U\jv) = 0$ for each $\vv \in E_i$.
\item The characteristic vector of $U$ is in 
$E_i^\perp =\bigoplus_{ j \ne i} E_j$. 
\end{enumerate}
\end{lemma}

\begin{lemma} \label{lem:mu_0_main}
Let $(V,E)$ be a $k$-regular graph. 
Let $V_1,\ldots, V_t$ be the connected components of $(V,E)$.  
The notation is defined as above. 
Let $U_i=U \cap V_i $. 
Then, the following are equivalent.
\begin{enumerate}
    \item $\mu_0$ is not a main eigenvalue of $2\D'-(a+b)\J$. 
    \item $|U_i|=|V_i|/2$ for each $i \in \{1,\ldots, t\}$. 
\end{enumerate} 
\end{lemma}

\begin{theorem} \label{thm:regular_angle}
Let $(V,E)$ be a $k$-regular graph of order $n$ with connected components $V=V_1\cup \cdots \cup V_t$. 
Let $\pi=\{U, V\setminus U\}$ be an equitable partition of $(V,E)$ with quotient matrix $\B$, where $U\ne \emptyset,V$. Let $U_i=U \cap V_i$. 
Let $\lambda_i$ be an eigenvalue of $(V,E)$, and $\mu_i$ that of  $2\D'-(a+b)\J$ given in Lemma~\ref{lem:eigen_2D'}. 
If $\lambda_0=k$, $\lambda_1$ are the eigenvalues of $\B$,   
then the following hold.
 \begin{enumerate}
\item If $|U_i|\ne |V_i|/2$ for some $i\in \{1,\ldots, t\}$, then the main eigenvalues of $2\D'-(a+b)\J$ are $\mu_0$, $\mu_1$. 
\item If $|U_i|= |V_i|/2$ for each $i \in \{1, \ldots, t\}$, then the main eigenvalue of $2 \D'-(a+b)\J$ is $\mu_1$. 
\end{enumerate}
\end{theorem}

\begin{proof}
Let $\boldsymbol{\pi}$ be the characteristic matrix of $\pi$.
For $i\ne 0$, if $\mu_i$ is a main eigenvalue of $2\D'-(a+b)\J$, then $\vv^\top (\I_U \jv) \ne 0$ for some $\vv \in E_i$ by Lemma \ref{lem:switch_main}. 
If $\vv^\top (\I_U \jv) \ne 0$ for some $\vv=(\vv_x)_{x \in V} \in E_i$, then 
\[
(\vv_x)^\top (\I_U \jv)= -\sum_{x \in U}\vv_x+\sum_{x \in V\setminus U} \vv_x \ne 0,
\]
and $\vv^\top  \boldsymbol{\pi}=(\sum_{x \in U} \vv_x,\sum_{x \in V \setminus U} \vv_x)$ is not the zero vector. 
Then, $\vv^\top  \boldsymbol{\pi}$ is a left eigenvector of $\B$, and hence $\lambda_i$ is an eigenvalue of $\B$. 
Therefore, for $i\ne 0$, the candidate of main eigenvalue $\mu_i$ is only $\mu_1$. 

If $|U_i|\ne |V_i|/2$ for some $i\in \{1,\ldots, t\}$, then $\mu_0$ is a main eigenvalue from Lemma \ref{lem:mu_0_main}. 
If $\mu_0$ is the only main eigenvalue, then $\jv \in \I_U E_0$.  
Since $\jv \not\in \I_U E_0={\rm Span}\{\I_U \jv_1, \ldots, \I_U \jv_t\}$, the eigenvalue $\mu_1$ must be main. 

If $|U_i|= |V_i|/2$ for each $i \in \{1, \ldots, t\}$, then $\mu_0$ is not a main eigenvalue from Lemma \ref{lem:mu_0_main}. The remaining candidate $\mu_1$ must be main. 
\end{proof}

\end{document}
